% PDF 197--201；原书第 168--172 页。
\section{教程——在资源受限时微调大语言模型}
\label{sec:ch04-tutorial}
\subsection{概述}
前面已经介绍几种参数高效微调技术，并概述了大语言模型微调的两种主要方法，即指令调优和对齐调优。本教程用 LoRA 和 QLoRA 训练大语言模型，使它完成一项基于指令的特定任务。这里集中在一项任务上，覆盖范围比严格意义上的指令调优要窄，但构造模板时采用指令调优的方法。

\begin{keypoint}
\noindent\textbf{目标}
\begin{itemize}
\item 演示参数高效微调在内存需求和输出质量方面的优点。
\item 比较较大的大语言模型和规模较小的大语言模型的能力。
\item 为生成的文本制定评分标准，使用能力更强的大语言模型评分。
\end{itemize}
\end{keypoint}

这里给出教程的简明版本。完整版见 \url{https://github.com/springer-llms-deep-dive/llms-deep-dive-tutorials}。

\subsection{实验设计}
本教程要构建一个大语言模型，输入客户与客服人员的对话，输出概括主要内容的摘要。这里记录的结果来自一个配有 16~GB V100 GPU 的 Google Colab 会话。使用的数据集是 TWEETSUMM（\refcite{ch04-feigenblat}{Feigenblat et al., 2021}），由客户和不同公司客服人员在 \texttt{x.com}（原 Twitter）上来回交流的对话组成。每段对话配有人工撰写的两句摘要，记录客户的请求和客服人员的回复。大多数对话有不同标注者编写的多份摘要。

为评估大语言模型所生成摘要的质量，按下列三项标准计分。

\begin{enumerate}
\item 对客户问题或投诉的描述是否基本准确？
\item 对客服人员回复的描述是否基本准确？
\item 摘要是否恰好由两句话组成？
\end{enumerate}

每满足一项标准，摘要得一分。依照 \refcite{ch04-dettmers-2023}{Dettmers 等人（2023）}的做法，用 GPT-4 给摘要评分。把含有上述标准的评分细则连同输入对话和生成的摘要交给 GPT-4，让它返回一个满分为 3 分的分数。

首先测试 DistilGPT-2。这是一个有 8500 万参数的自回归大语言模型，由 GPT-2 提供监督来训练。选择它是因为内存需求相对较低，便于在本实验的 Colab 环境中微调。

随后尝试改用更大的大语言模型来改善结果。更丰富的语言知识有助于模型理解这些消息中发生了什么。为此，采用 Llama-2-7B，这是 Meta 于 2023 年发布的、具有 70 亿参数的自回归文本生成模型。虽然它的能力更强，但用微调 DistilGPT-2 的方法尝试微调它时，显存不足。这就需要采用参数高效微调技术。接下来使用 LoRA 和 QLoRA，比较不同训练方法得到的模型表现与训练时间。

\subsection{结果与分析}
\subsubsection{DistilGPT-2}
作为参照，先让 DistilGPT-2 在不做微调的情况下为测试集中的每段对话生成摘要。为文本生成定义一个 \texttt{transformers} 流水线，然后输入由 TWEETSUMM 测试集按模板构造的提示。输出质量较差。这并不意外，因为 DistilGPT-2 的规模较小，未经进一步微调，难以展现突出的涌现能力。接着用 Python 包 \texttt{trl} 在训练数据上微调模型。它对 \texttt{transformers} 的功能作了便于使用的封装。微调后的 DistilGPT-2 表现优于基础模型，尤其是在摘要长度这一项上，但对客户和客服人员对话的描述仍然不够准确。

为考察总体表现，分别用基础模型和微调模型为测试集中的 50 段对话生成摘要，再用 GPT-4 评分。基础模型摘要的累计得分为 2 分，满分为 150 分。这个结果很差，无法满足任务要求。微调模型的表现好得多，得到 67/150 分，不过仍与理想效果相差较远。

\begin{table}[!htbp]
\centering
\caption{在 TWEETSUMM 训练集上用各方法调优模型，并在测试集上推理得到的最终分数，满分为 150 分。表中同时列出各模型的调优时间。}
\label{tab:ch04-tweetsumm}
\small
\begin{tabularx}{\linewidth}{@{}Xrr@{}}
\toprule
模型配置 & 摘要得分（满分 150） & 调优时间（分钟）\\
\midrule
基础 DistilGPT-2 & 2 & 0\\
全参数微调的 DistilGPT-2 & 67 & 9.7\\
LoRA 调优的 DistilGPT-2 & 58 & 6.9\\
QLoRA 调优的 DistilGPT-2 & 52 & 14.3\\
基础 Llama-2-7B & 25.5 & 0\\
全参数微调的 Llama-2-7B & \multicolumn{2}{c@{}}{失败}\\
LoRA 调优的 Llama-2-7B & 131 & 75.1\\
QLoRA 调优的 Llama-2-7B & 125 & 21.3\\
\bottomrule
\end{tabularx}
\end{table}

如第~\ref{sec:ch04-lora}~节所述，低秩适配器是一种常用且高效的降低训练内存需求的方法。只需调优两个低秩矩阵，推理时将所得的权重更新加到原权重中，无需微调整个权重矩阵。这大幅减少了训练时需要在内存中保存梯度的参数数量。我们还试用效率更高的 QLoRA，先把模型权重量化到 4 位，再用 LoRA 方法调优。

LoRA 和 QLoRA 在 TWEETSUMM 数据集上的比较见表~\ref{tab:ch04-tweetsumm}。二者的表现都没有达到全参数微调的水平，但远好于基础模型。尽管 DistilGPT-2 的表现较低，这两种方法训练时的总 GPU 负担较小。与全参数微调相比，LoRA 调优的峰值 GPU 显存占用少 228~MB，QLoRA 调优少 336~MB。考虑到 DistilGPT-2 的权重矩阵约占 356~MB，这一差别相当明显。

\subsubsection{Llama-2-7B}
接着尝试用更大的语言模型来改善结果。它拥有更丰富的语言知识，可能更有能力理解这些消息中发生了什么。Llama-2-7B 符合要求。重复基础模型的零样本摘要实验，Llama-2 得到 25.5/150 分。表现依然较差，但较基础 DistilGPT-2 明显改善。接着尝试全参数微调 Llama-2，GPU 显存不足。Llama-2-7B 有 70 亿参数，仅模型权重便占约 12~GB 显存，加入微调所需的梯度后，总占用增至约 64~GB，远超本实验所用 V100 GPU 的 16~GB 显存。

\sourcefigure{fig04_12.pdf}{各种推理配置在 TWEETSUMM 摘要任务上的最终得分。上图显示原始分数，满分为 150 分。下图按三项标准拆分得分，即是否准确概括客户的发言、是否准确概括客服人员的发言，以及输出是否恰好两句话。由于显存不足，Llama-2-7B 的全参数微调未能产出模型。}{fig:ch04-tweetsumm}

因此，在 Llama-2-7B 上改用 PEFT 方法，训练得以开始，未出现 CUDA 错误。调优一轮，LoRA 用时 75 分钟，QLoRA 只用 21 分钟。这样做使表现大幅改善。测试集上，LoRA 调优模型得 131/150 分，QLoRA 调优模型得 125 分。%
\footnote{译注：原书正文将 QLoRA 的得分印为“125/131”。表~\ref{tab:ch04-tweetsumm} 和图~\ref{fig:ch04-tweetsumm} 均以 150 为满分，列出的 QLoRA 得分为 125。}

图~\ref{fig:ch04-tweetsumm} 汇总本教程各配置在测试集上的评估结果。两个经过适配器调优的 Llama-2-7B 模型总分最高，在每项评分标准上也表现最好。从下图可见，微调后的 DistilGPT-2 学会了将摘要限制为两句，但摘要仍不够准确，未达到 GPT-4 的评分要求。基础 Llama-2-7B 生成的摘要中，被判为准确的数量与全参数微调的 DistilGPT-2 相同，但没有经过进一步训练，就无法遵守格式要求。

\subsection{结论}
实验表明，较小的大语言模型可通过调优学会遵循特定指令，但受信息容量限制，其语义处理能力最终难与更大的模型相比。在两个经过调优的 Llama-2 模型中，QLoRA 略逊于 LoRA，但其调优时间不到后者的三分之一。对训练数据量较大的场景，这项权衡十分重要。总体而言，低秩适配器调优利用了 Llama-2-7B 的大量参数，得到了能够准确概括对话内容的摘要机器人。
